\documentclass[11pt]{article}
\usepackage[a4paper,margin=27mm]{geometry}
\usepackage{amsmath,amssymb,amsthm,mathtools,microtype}
\usepackage[colorlinks=true,allcolors=blue]{hyperref}
\newtheorem{theorem}{Theorem}
\newtheorem{lemma}[theorem]{Lemma}
\newtheorem{proposition}[theorem]{Proposition}
\newtheorem{corollary}[theorem]{Corollary}
\newcommand{\C}{\mathbb C}
\newcommand{\E}{\mathbb E}
\newcommand{\Tr}{\operatorname{Tr}}
\newcommand{\I}{\mathcal I}
\newcommand{\CN}{\mathcal{CN}}
\newcommand{\capq}{C_{\mathrm{opt}}}
\title{Classical limits and robust Gaussian coding gaps\newline
for compound bosonic fading channels}
\author{Lluis Eriksson}
\date{26 October 2026}
\begin{document}
\maketitle
\begin{abstract}
We study unassisted classical communication through memoryless mixtures
of noisy, phase-insensitive bosonic Gaussian channels. Both thermal noise
and photon budget scale with an energy parameter. An explicit
measure--prepare representation converts every quantum block code,
including entangled codewords and collective receivers, into a classical
code with identical error and an exactly quantified cost increment.
The same simulation works for an unknown fixed mixing law.
A Gaussian-smoothing estimate and a uniform finite-alphabet reduction
give an explicit \(O(E^{-1})\) comparison with a classical compound
variational value. This value optimizes one input distribution before
taking the worst mixing law. We also prove a strict separation
between classical optimal modulation and the time-sharing hull of Gaussian
coherent modulations for a matched-reset example. Matching bounds determine the high-energy limit of a regularized
Holevo functional with Gaussian average marginals. A block-uniform,
energy-controlled continuity modulus removes both exact marginal
Gaussianity and Gaussianity of the full average state. We prove that
almost every product coherent Gaussian random codebook from a fixed
finite deterministic schedule has the required approximate marginals,
controlling empirical moments separately. Strict cost slack or an
explicit common rescaling enforces physical feasibility. This yields
an operational converse for reliable sequences in the defined typical
class, below an unrestricted non-Gaussian common-code rate at sufficiently
large energy. It does not assert a good decoder for every random book
or cover every protocol described as Gaussian. The Gaussian, coding and
continuity tools used in the argument have established antecedents.
\end{abstract}

\section{Model and operational statement}

All logarithms are natural, and \(h_2\) denotes binary Shannon entropy.

Write
\(\phi_v(z)=(\pi v)^{-1}e^{-|z|^2/v}\), so that
\(\E|G_v|^2=v\). A thermal state \(\gamma_t\) has mean photon
number \(t\). Let \(\Phi_{a,t}\) be the phase-insensitive channel
sending a coherent input \(|z\rangle\) to
\(D(az)\gamma_tD(az)^*\). Its characteristic function is
\[
 \chi_{\Phi_{a,t}(\rho)}(\xi)=
 \chi_\rho(\bar a\xi)
 e^{-[t+(1-|a|^2)/2]|\xi|^2}.
\]
Physicality requires \(t\geq \max(0,|a|^2-1)\).

Fix \(K<\infty\), \(0<v_*\leq v^*<\infty\), and the compact parameter
space
\[
 \Theta=\{(a,v)\in\C\times\mathbb R: |a|^2\leq K,\quad
                                    v_*\leq v\leq v^*\}.
\]
Let \(\Gamma\) be a nonempty weakly closed set of probability laws on
\(\Theta\). For \(\nu\in\Gamma\) define
\[
 N_{E,\nu}=\int_\Theta\Phi_{a,vE}\,d\nu(a,v),\qquad
 W_\nu(y|x)=\int_\Theta\phi_v(y-ax)\,d\nu(a,v).
\]
The unknown law \(\nu\) is fixed throughout each code; conditional on it,
branches are sampled independently in different uses. Neither party
receives the branch label or the unknown law. This is a compound
memoryless model, not an arbitrarily varying channel.

A quantum code with \(M\) equally likely messages has \(m\)-mode states
\(\rho_j\), a single decoding POVM \(\{D_j\}\), and expected photon cost
\[
 M^{-1}\sum_j\Tr\rho_j\sum_{i=1}^m n_i\leq mE .
\]
The average message error must tend to zero uniformly over \(\nu\in\Gamma\).
Write \(\capq(E,\Gamma)\) for its supremum of rates. It is an
unassisted \emph{classical} capacity of a quantum channel, not a
capacity for transmitting qubits.

For a law \(\mu\) on \(\C\), write \(I_\mu(W_\nu)=I(X;Y)\), and set
\[
 V_\Gamma(P)=
 \sup_{\mu:\int|x|^2d\mu\leq P}\inf_{\nu\in\Gamma} I_\mu(W_\nu).
\]
It is essential that the input optimization precede the worst-law
operation. This is generally different from the worst separately
optimized capacity.

\begin{theorem}[Uniform classical comparison]\label{thm:main}
For \(E>0\), \(E\geq 2K/v_*\), the preceding model obeys
\[
 V_\Gamma(1)-\frac{2}{v_*E}
 \ \leq\ \capq(E,\Gamma)\
 \leq V_\Gamma(1)+
 \frac{\log(1+K/v_*)+4K/v_*}{E}.                 \tag{1}
\]
For \(K=0\) both operational capacities are zero.
The constants are independent of \(\Gamma\); \(K,v_*,v^*\) stay fixed.
\end{theorem}

Coding length tends to infinity at fixed \(E\), before energy tends to
infinity. No uniform blocklength, efficient optimizer, fixed-error strong
converse or positivity for every family is asserted. In particular,
a family containing a completely replacing channel has value zero.
If \(\Gamma=\Gamma_E\) changes between energy levels, (1) compares to
\(V_{\Gamma_E}(1)\); it does not assert that this sequence converges.

\section{An exact simulation preserving the common decoder}

Put \(k=|a|^2\). For \(t\geq k\), the heterodyne measure--prepare
representation is
\[
 \Phi_{a,t}(\rho)=\int_\C Q_\rho(\alpha)
           D(a\alpha)\gamma_{t-k}D(a\alpha)^*\,d^2\alpha,\qquad
 Q_\rho(\alpha)=\pi^{-1}\langle\alpha|\rho|\alpha\rangle.       \tag{2}
\]
This is a specialization of the established Gaussian
entanglement-breaking representation \cite{holevoEB,holevoReview}.
Indeed, the characteristic function of the heterodyne outcome is
\(\chi_\rho(\eta)e^{-|\eta|^2/2}\); multiplying it by the preparation
factor \(e^{-(t-k+1/2)|\xi|^2}\), at \(\eta=\bar a\xi\), gives
the characteristic function in Section 1. Thus the noise subtraction
in (2) is \(k\), and
\[
 \int|\alpha|^2Q_\rho(\alpha)\,d^2\alpha=\Tr\rho n+1.          \tag{3}
\]
The identity extends to correlated \(m\)-mode Husimi outcomes, with
total cost \(\Tr\rho\sum_i n_i+m\).

For \(t>k\), refine the prepared thermal state by its positive
coherent-state representation. For every quantum code this gives the
classical Markov chain
\[
 J\longrightarrow A^m\longrightarrow B^m\longrightarrow\widehat J,
 \quad
 P(\widehat J=j|B^m=\beta^m)
   =\langle\beta^m|D_j|\beta^m\rangle .                       \tag{4}
\]
Given \(A^m\), the middle transition is a product of mixtures with
variance \(vE-k\). Both the stochastic Husimi encoder and the final Born
decision kernel are independent of \(\nu\). Equation (4) reproduces
the error of the quantum code for \emph{every} law, with its correlations
intact. The hidden \(B^m\) is a converse simulation variable, not an
observation supplied to the physical receiver.

Normalize \(X=A/\sqrt E\), \(Y=B/\sqrt E\). Define
\[
 W_{E,-,\nu}(y|x)=\int\phi_{v-k/E}(y-ax)\,d\nu,\qquad
 W_{E,+,\nu}(y|x)=\int\phi_{v+1/E}(y-ax)\,d\nu .
\]
The lower channel is realized by coherent input \(|\sqrt E x\rangle\)
and heterodyne reception. The upper simulation has mean cost at most
\(1+1/E\), by (3). Consequently the operational classical compound
capacities satisfy the exact sandwich
\[
 C_{\rm cmp}(W_{E,+},\Gamma,1)
 \leq \capq(E,\Gamma)
 \leq C_{\rm cmp}(W_{E,-},\Gamma,1+1/E).                     \tag{5}
\]
No additivity or Gaussian optimizer of the quantum mixture is used.
For continuous \(\nu\), the state integrals are trace-norm Bochner
integrals; probabilities and nonnegative moments also permit Tonelli.

\section{Uniform information estimates}

\begin{lemma}[Entropy continuity]\label{lem:entropy}
Suppose planar densities \(f_n,f\leq L<\infty\) converge in \(L^1\)
and their second moments are at most \(M<\infty\). Then
\(h(f_n)\to h(f)\).
\end{lemma}
\begin{proof}
On a fixed ball use the uniform continuity of \(-u\log u\) on \([0,L]\).
Outside radius \(R\), let \(\delta\) be the mass. The restricted,
normalized density has second moment at most \(M/\delta\). Gaussian
maximum entropy therefore gives
\[
 -\delta\log L\leq-\int_{|y|>R} f\log f
 \leq\delta\log(\pi eM/\delta^2).
\]
Here \(\delta\leq M/R^2\); both bounds tend uniformly to zero as
\(R\to\infty\), with value zero when \(\delta=0\).
This also proves finiteness of the entropies in the assertion.
\end{proof}

\begin{lemma}[Noise sensitivity]\label{lem:noise}
Fix any input law of finite second moment and any mixture of bounded
gains. If component variances follow
\(v_\theta=v_0+\theta d\in[\ell,U]\), \(\ell>0\), with
\(|d|\leq D\), then
\[
 |I_\mu(W_{\theta_1})-I_\mu(W_{\theta_0})|
       \leq \frac{2D}{\ell}|\theta_1-\theta_0|.              \tag{6}
\]
The constant is independent of the input peak, its cost and the mixing law.
\end{lemma}
\begin{proof}
First take a compactly supported input. Let
\(f_{a,v}=(a_\#\mu)*\phi_v\), \(f=\int f_{a,v}d\nu\), and
\(J=\int d\,\nabla f_{a,v}d\nu\). Complex variance \(v\) means
real-coordinate variance \(v/2\), hence
\(\partial_v\phi_v=\Delta\phi_v/4\).
Differentiation and integration by parts yield
\[
 \partial_\theta h(f)=\tfrac14\int J\cdot\nabla\log f .
\]
The Gaussian score conditional-expectation identity and Jensen give
\(F(f_{a,v})\leq4/v\leq4/\ell\), where
\(F(f)=\int |\nabla f|^2/f\) is planar Fisher information.
Weighted Cauchy--Schwarz gives
\[
 F(f)\leq4/\ell,\qquad \int |J|^2/f\leq4D^2/\ell .
\]
Thus \(|\partial_\theta h(f)|\leq D/\ell\).
To justify the differentiations, bounded centers and variances give
common Gaussian upper tails, a Gaussian lower density bound, and
polynomial-times-Gaussian bounds for derivatives; boundary terms vanish.
These bounds hold uniformly over atoms or continuous mixtures.

For a general input, replace its tail outside radius \(R\) by zero.
The coupled outputs converge in \(L^1\), are bounded by
\((\pi\ell)^{-1}\), and have second moment at most
\(K\E|X|^2+U\). Lemma~\ref{lem:entropy} passes the integrated entropy
inequality to the limit. This does not require differentiability for the
untruncated heavy-tailed input.
For a fixed letter \(x\), the same entropy bound is independent of \(x\).
Conditional entropies are integrable, bounded below by a constant and
above by \(\log[\pi e(K|x|^2+U)]\). Integrate the conditional bound and
subtract it from the output bound to obtain (6).
\end{proof}

Revealing the independent branch only for an information upper bound gives
\[
 0\leq I_\mu(W_\nu)
 \leq\int\log(1+|a|^2P/v)\,d\nu
 \leq\log(1+KP/v_*).                                      \tag{7}
\]
No branch is revealed to the actual decoder. The function \(V_\Gamma\)
is finite, nondecreasing and concave, and \(V_\Gamma(0)=0\).
Indeed input concavity holds simultaneously for every law and survives
the infimum. In particular
\[
 V_\Gamma(1+u)\leq(1+u)V_\Gamma(1).                         \tag{8}
\]

\section{One code for a continuous compound family}

We supply the reduction to the established finite-alphabet compound
coding mechanism \cite{types,compound}; a finite-output theorem alone
does not justify continuous observations or physical costs.

\begin{proposition}\label{prop:compound}
For the Gaussian-smoothed family above and \(P>0\),
\[
 C_{\rm cmp}(W,\Gamma,P)=V_\Gamma(P).                        \tag{9}
\]
The same holds for either family in (5) at fixed admissible \(E\).
\end{proposition}
\begin{proof}
For the converse, Fano and data processing applied to a common stochastic
encoder give, for each \(\nu\),
\[
 (1-\epsilon)\log M-h_2(\epsilon)
 \leq I(J;Y_\nu^m)\leq\sum_i I_{\mu_i}(W_\nu)
 \leq mI_{\bar\mu}(W_\nu),\qquad
 \bar\mu=m^{-1}\sum_i\mu_i .
\]
The law \(\bar\mu\) has cost at most \(P\) and is independent of \(\nu\).
Taking the infimum over \(\nu\), then optimizing feasible \(\bar\mu\),
gives the required order. The block inequality can be proved without
subtracting infinite differential entropies: the KL chain rule with
product output marginals gives
\[
 \sum_i I(X_i;Y_i)=I(X^m;Y^m)
             +D(P_{Y^m}\Vert\textstyle\prod_i P_{Y_i}).
\]
All terms are finite in our bounded-noise regime by (7).

For achievability choose a feasible input near \(V_\Gamma(P)\).
Mixing it with a small mass at zero makes its cost strictly below \(P\);
concavity loses arbitrarily little information uniformly in \(\nu\)
by (7). Truncate its tail to zero, then quantize a bounded input ball.
The Gaussian kernel has the global input-gradient bound
\(\sqrt{2K}/(\pi\sqrt e\,v_*^{3/2})\).
Thus their density difference is at most this constant times
\(\E|X_l-X|\). Under this coupling the output densities converge uniformly
on every output ball, uniformly in \(\nu\), and have common density and
second-moment bounds. The proof of Lemma~\ref{lem:entropy}, split into
that ball and its complement, gives uniform output entropy convergence.
The conditional entropy is jointly continuous on compact input and
parameter domains and is bounded by a constant plus
\(\log(1+|x|^2)\). Its tail integral under the fixed input vanishes.
Thus one finite input alphabet with strict cost slack approximates all
the information quantities simultaneously.

For this finite input law choose nested finite output partitions:
expanding dyadic boxes, their refined cells, and one exterior cell.
They generate the Borel sigma algebra. Their mutual informations
increase to the continuous-output information by posterior martingale
convergence; the posterior entropy is bounded by the finite input entropy.
Each partition information is continuous in \(\nu\): Gaussian kernels
depend continuously on \((a,v)\) in \(L^1\), so every Borel cell
probability is continuous, including the unbounded exterior cell.
The limit is also continuous: bounded means and compact positive
variances give common Gaussian tails and local density convergence.
Dini's theorem on compact \(\Gamma\) now supplies one finite partition
with uniformly small loss. This partition is fixed before code length.

For completeness, let \(A,B\) be these finite alphabets, and let
\(R<\inf_W I(p,W)\). Choose constant-composition types \(p_n\to p\)
on its positive support, retaining strict cost slack.
Draw \(M=\lceil e^{nR}\rceil\) codewords uniformly and independently
from their type class. Decode a unique codeword whose empirical mutual
information with the output exceeds \(R+\eta\), with \(\eta>0\)
below the rate slack. This decoder is independent of \(W\).
Hoeffding bounds for output counts within each input-symbol class,
followed by finite-alphabet entropy continuity, make the transmitted
word fail the threshold with probability \(e^{-cn}\), uniformly in \(W\).
For any fixed output word a random competing codeword crosses the
threshold with probability at most
\(\operatorname{poly}(n)e^{-n(R+\eta)}\):
a conditional type class has size at most \(e^{nH(X|Y)}\),
whereas the input type class has size at least
\((n+1)^{-|A|}e^{nH(p_n)}\). Sum over the polynomial number of joint
types and over competitors.

To choose one deterministic codebook, cover the finite stochastic-matrix
family in per-row total variation by a net of radius \(n^{-2}\),
of cardinality at most \((Cn^2)^{|A||B|}\).
Choose its points from the family; all retain the rate slack.
The expected \emph{sum} of average errors on this net tends to zero,
so a deterministic codebook has uniformly vanishing net error.
Product-channel total variation changes every decoder error by at most
\(n n^{-2}\). Therefore its error vanishes over the whole family.
There is no exchange of an expected supremum with a supremum of expectations.

Each word has exact additive cost \(n\sum_xp_n(x)|x|^2<nP\).
Mapping \(x\) to \(|\sqrt E x\rangle\) enforces physical expected
photon cost \(E\sum_i|x_i|^2\) for every message.
Fix the slack and finite approximations first, let \(n\to\infty\),
then remove those approximations. Equivalently, for each rate strictly
below the variational value, choose all finite approximations once before
taking the coding limit. If that value is zero, the converse already
settles the zero-rate case. The result asserts no uniform
coding length or efficient search.
\end{proof}

\begin{proof}[Proof of Theorem~\ref{thm:main}]
For the plus family use (6) with \(D=1,\ell=v_*\), over an interval
of length \(1/E\). For the minus family use \(D=K,\ell=v_*/2\),
valid when \(E\geq2K/v_*\). The bounds survive infima and suprema over
the same laws and input constraints. Consequently
\[
 V_{E,+,\Gamma}(1)\geq V_\Gamma(1)-2/(v_*E),\qquad
 V_{E,-,\Gamma}(P)\leq V_\Gamma(P)+4K/(v_*E).
\]
Apply (9) in (5), then (8) and (7). This gives (1).
The optimized coherent-modulation heterodyne rate is within
\([\log(1+K/v_*)+(4K+2)/v_*]/E\) of the unrestricted quantum rate.
\end{proof}

\section{A matched-reset application: a strict Gaussian-hull gap}

Fix \(k,r>0\), \(0<\lambda<1\), \(q=1-\lambda\), \(s=k+r\), and
\[
 W_\theta(y|x)=\lambda\phi_r(y-\sqrt k e^{i\theta}x)+q\phi_s(y).
\]
For \(\mu_0=\CN(0,1)\), the average output is exactly \(\phi_s\).
Let \(F\) denote the signal/reset flag and
\(H(x)=H(F|Y,X=x)\). It is radial, bounded, positive at finite \(x\),
and tends to zero at infinity. The last assertion follows by separating
the two widely separated Gaussian clouds and applying binary Fano.
It follows that \(H\) is not affine in \(|x|^2\).

There exists a compactly supported bounded radial \(h\), normalized by
\(\|h\|_\infty\leq1\), with
\[
 \int h\,d\mu_0=\int |x|^2h\,d\mu_0=0,\qquad
 a_0=\int hH\,d\mu_0>0.                                   \tag{10}
\]
Indeed, if all such moment-annihilating tests annihilated \(H\), two
fixed radial moment correctors would imply that \(H\) equals an affine
function of \(|x|^2\) on every bounded radial domain, hence globally
almost everywhere. Boundedness and decay would make it zero,
contradicting positivity. Reverse the sign and normalize a nonzero test.
The moment-corrector matrix is invertible by taking two disjoint annuli
with distinct conditional squared radii.

For \(\mu_\epsilon=(1+\epsilon h)\mu_0\), (10) preserves mass and cost.
Since the output second moment stays \(s\), Gaussian cross entropy and
the binary flag identity give
\[
 I_{\mu_\epsilon}(W_0)-I_{\mu_0}(W_0)
  =\epsilon a_0-D(p_{Y,\epsilon}\Vert\phi_s)
  \geq\epsilon a_0-\epsilon^2\int h^2\,d\mu_0.               \tag{11}
\]
The last inequality is KL data processing and
\(\log(1+u)\leq u\). A fixed sufficiently small \(\epsilon>0\)
gives a positive, non-numerical improvement.

Let \(G_\vartheta\) range over all real two-dimensional Gaussian laws
on \(\C\), including displaced, anisotropic and singular laws.
Write \(\vartheta=(m,V)\), \(p(\vartheta)=|m|^2+\Tr V\),
\(J(\vartheta)=I_{G_\vartheta}(W_0)\), and define
\[
 H_G=\sup_{\eta:\int p\,d\eta\leq1}\int J\,d\eta .           \tag{12}
\]
Here the time-sharing setting is known to both parties and does not
itself carry a message.

\begin{proposition}[Strict Gaussian scheduling gap]\label{prop:hull}
For this matched-reset channel, \(H_G<C(1)\), where
\(C(P)=\sup_{\mu:\E|X|^2\leq P}I_\mu(W_0)\).
\end{proposition}
\begin{proof}
The symmetry, injectivity and strict-growth strategy has nearby
antecedents in fading problems \cite{catalano,tu}; we check it directly
for the present channel and Gaussian scheduling class.
The output-law map is injective: subtract \(q\phi_s\) and Fourier
transform the remaining convolution, whose Gaussian factor is nowhere
zero. Thus output entropy, and hence mutual information, is strictly
concave in the input law.

If a Gaussian law attained \(C(1)\), rotation covariance and strict
concavity would force it to be circular \(\CN(0,P)\).
The finite concave nondecreasing capacity \(C(P)\) is unbounded, since
revealing the independent flag costs at most \(h_2(\lambda)\) and gives
\[
 C(P)\geq\lambda\log(1+kP/r)-h_2(\lambda).
\]
Such a function is strictly increasing: a flat interval would, by
concavity and monotonicity, force all subsequent slopes to vanish.
Therefore \(P=1\), contradicted by (11). No Gaussian law attains \(C(1)\).

The parameter sublevel sets of \(p\) are compact. Laws \(\eta\) with
\(\int p\,d\eta\leq1\) are tight and weakly closed, hence compact.
Gaussian smoothing and Lemma~\ref{lem:entropy} show that \(J\) is continuous,
including at singular covariance matrices. Furthermore
\[
 0\leq J(\vartheta)\leq\lambda\log(1+kp(\vartheta)/r).
\]
The upper bound is sublinear in \(p\), so its tails over \(p>R\)
are uniformly integrable under the mean-cost constraint.
Compactly supported continuous truncation then makes
\(\eta\mapsto\int J\,d\eta\) continuous. The supremum (12) is attained.

For an attaining \(\eta_*\), the barycenter
\(\bar\mu=\int G_\vartheta\,d\eta_*\) has cost at most one.
To justify averaging entropies, every output density is bounded below by
\(q\phi_s\) and above by a common constant. Thus
\(|\log p_{Y,\bar\mu}(y)|\leq C+C'|y|^2\), and the average output
second moment is finite. Individual output entropies have integrable
bounds by constants and by \(C+\log(1+p(\vartheta))\); the conditional
entropy of \(W(.|x)\) is bounded through its binary flag identity.
Fubini therefore gives the exact identity
\[
 I_{\bar\mu}(W_0)-H_G
   =\int D(p_{Y,\vartheta}\Vert p_{Y,\bar\mu})\,d\eta_*\geq0 .
\]
If \(H_G=C(1)\), equality forces almost all output laws to coincide.
Injectivity makes almost all Gaussian input laws identical to
\(\bar\mu\), which would be a Gaussian maximizer. This contradiction
proves the strict gap. Its size and its optimizer are not computed.
\end{proof}

For any correlated multivariate Gaussian modulation on product coherent
states with total cost \(mE\), positive thermal preparation gives exactly
\[
 \chi(X^m:N_{E,\theta}^{\otimes m})
 \leq I(X^m:Y^m)\leq\sum_i I_{\mu_i}(W_\theta)\leq mH_G.     \tag{13}
\]
The first inequality uses the preparation chain
\(X^m\to Y^m\to|\sqrt E Y^m\rangle\), valid for coherent letters
without the input heterodyne cost increment.
Here \(N_{E,\theta}=\lambda\Phi_{\sqrt k e^{i\theta},rE}
+q\gamma_{sE}\Tr\); we may take \(E\geq2k/r\), which also
ensures physicality and the hypotheses of Theorem~\ref{thm:main}.
The last uses Gaussian marginals and their average cost, allowing
individual costs above one. Known, non-message-bearing time sharing of
such blocks is also covered. This does \emph{not} bound squeezed or
entangled Gaussian quantum signal states, message-bearing covariance
selection, or every protocol described as Gaussian.

Choose a finite input law of strict cost below one whose information
exceeds \(H_G\) by a fixed amount, using Proposition~\ref{prop:compound}'s
approximation argument. Heterodyne reception loses at most \(2/(rE)\)
by (6), so it eventually beats the entire benchmark (13).
For an unknown fixed phase,
\(W_\theta(y|x)=W_0(e^{-i\theta}y|x)\) implies
\(I_\mu(W_\theta)=I_\mu(W_0)\) for every input law.
The common code, rather than this equality alone, is supplied by (9).
No phase oracle is assumed.

\subsection{Ensembles with Gaussian average quantum states}

Consider any ensemble \(\{p_j,\rho_j\}\) on \(m\) modes with average
state \(\bar\rho=\sum_jp_j\rho_j\) Gaussian and
\(\Tr\bar\rho\sum_i n_i\leq mE\).
The average state may be displaced, squeezed or entangled across uses;
the conditional states need not be Gaussian. Continuous ensembles use
the relative-entropy definition of Holevo information. Then
\[
 \frac1m\chi\bigl(\{p_j,N_{E,\theta}^{\otimes m}(\rho_j)\}\bigr)
 \leq H_G+\frac{\lambda\log(1+k/r)+4k/r}{E}.                \tag{14}
\]
The simulation in (4) implies
\[
 \chi(J:\text{output})\leq I(J;B^m)
 \leq I(A^m;B^m)\leq\sum_i I_{\mu_i}(W_{E,-,\theta}).
\]
For continuous labels, the first inequality is
\(D(\sigma_j\Vert\bar\sigma)\leq D(P_{B|j}\Vert P_B)\), integrated
in \(j\), without subtracting an entropy of the label. To justify it,
approximate the density ratios by conditional expectations on increasing
finite partitions of \(B\), preparing the average coherent state in each
cell. Finite-partition data processing, \(L^1\) convergence of the
ratios, trace-norm convergence of the prepared states, and lower
semicontinuity of relative entropy pass to the limit. Tonelli averages
the inequalities. The remaining steps are classical data processing and
the block KL chain rule. Linearity of the Husimi function makes its average law
\(Q_{\bar\rho}\), which is multivariate Gaussian for any Gaussian
quantum state: heterodyne is a Gaussian measurement
\cite{gaussianReview}. Its normalized marginals \(\mu_i\) are real planar
Gaussian laws, with average cost at most \(1+1/E\).
Correlations in \(A^m\) are allowed in the displayed block inequality.
The noise-sensitivity bound, uniform in each marginal cost, gives
\[
 I_{\mu_i}(W_{E,-,\theta})\leq
 I_{\mu_i}(W_0)+\frac{4k}{rE}.
\]
Writing \(H_G(P)\) for (12) at cost \(P\), the sum is therefore
at most \(mH_G(1+1/E)+4mk/(rE)\).
The hull is finite, nondecreasing and concave in cost, with
\(H_G(0)=0\); mix its parameter probability laws to see concavity.
Thus \(H_G(1+1/E)\leq H_G+H_G/E\).
Revealing the independent signal flag and applying Jensen in cost
also yields \(H_G\leq\lambda\log(1+k/r)\), proving (14).
All output entropies are finite under the photon constraint and the
uniform preparation-noise floor. Known time sharing among such blocks
is covered by summing their marginal costs before applying \(H_G(P)\).

By Proposition~\ref{prop:hull}, a fixed finite strict-cost law
can have \(I_\mu(W_0)\geq H_G+\delta\) for some \(\delta>0\).
Its coherent-modulation heterodyne rate exceeds the supremum of (14),
over all block lengths, by at least
\[
 \delta-\frac{\lambda\log(1+k/r)+(4k+2)/r}{E},              \tag{15}
\]
which is positive for sufficiently large \(E\).
The common unknown-phase decoder is again supplied by (9).

Equation (14) uses the global Gaussian-average condition. Its
distinction from a Gaussian ensemble is standard; see
\cite[Remark 2]{holevoMeasurement}, where a thermal average of Fock
states gives an example. That work concerns Gaussian measurement
channels, not the present non-Gaussian signal/reset mixture.
Finite codebooks generally do not have an exactly Gaussian average.
Section~\ref{sec:robust} gives a quantitative bound under approximately
Gaussian average marginals and proves that a precisely defined typical
Gaussian random-codebook class satisfies it. Message-dependent covariance
choices and arbitrary deterministic books are not automatically covered.

\subsection{A matching limit for the ensemble benchmark}

Define the regularized Gaussian-average Holevo functional
\[
 R_G(E)=\sup_{m\geq1}\frac1m
 \sup_{\substack{\{p_j,\rho_j\}:\ \bar\rho\ \mathrm{Gaussian}\\
                   \Tr\bar\rho\sum_i n_i\leq mE}}
 \chi\bigl(\{p_j,N_{E,0}^{\otimes m}(\rho_j)\}\bigr).        \tag{16}
\]
\begin{corollary}[Gaussian-average Holevo limit]\label{cor:gausslimit}
For \(E\geq2k/r\), the functional (16) satisfies
\[
 H_G-\frac{2}{rE}\leq R_G(E)
 \leq H_G+\frac{\lambda\log(1+k/r)+4k/r}{E}.                \tag{17}
\]
In particular \(R_G(E)\to H_G\).
\end{corollary}
\begin{proof}
The upper bound is (14). For the lower, take a parameter probability law
nearly optimal in (12). Mixing it with a small mass at the zero-cost
point Gaussian creates strict cost slack with arbitrarily small loss
of its average \(J\). Replace its high-cost tail by that point. The
information loss tends to zero since
\(J(\vartheta)\leq\lambda\log(1+kp(\vartheta)/r)=o(p(\vartheta))\).
On a compact parameter sublevel, continuity of \(p,J\) permits finite
quantization; rational approximation of the resulting weights retains
the cost slack. Thus rational finite time-sharing laws approximate
\(H_G\) from below with strict mean cost below one.

For weights with denominator \(m\), use \(m\) independent coherent
Gaussian modulations in their respective parameter settings. Their
average quantum input is a product Gaussian state, of photon cost
\(E\sum_i p(\vartheta_i)<mE\). Heterodyne information adds over these
independent modes, and Holevo information is at least the information
of that physical measurement. Lemma~\ref{lem:noise} gives at least
\(\sum_i J(\vartheta_i)-2m/(rE)\). Divide by \(m\) and remove the
finite approximations at fixed \(E\). This proves the lower bound
in (17) without claiming a finite-message code with exactly Gaussian
average.
\end{proof}

The quantum channels satisfy
\(N_{E,\theta}=\mathcal U_\theta\circ N_{E,0}\), where
\(\mathcal U_\theta\) rotates the output and leaves the thermal reset
invariant. Consequently (16) is phase independent even for arbitrary
conditional input states. For the compact family of unknown fixed
phases, the classical compound value is \(C(1)\), by the output-rotation
identity. Equations (1), (17) and Proposition~\ref{prop:hull} give
the exact limiting separation
\[
 \lim_{E\to\infty}
 \bigl[\capq(E,\text{all fixed phases})-R_G(E)\bigr]
       =C(1)-H_G>0.                                      \tag{18}
\]
The common decoder in the operational term comes from (9), whereas
\(R_G\) remains the explicitly defined Holevo functional. Its block
supremum is not silently replaced by an operational constrained-code
capacity. The positive constant in (18) is not numerically evaluated.

\section{Robust marginal bounds and typical Gaussian codebooks}
\label{sec:robust}

The full-block Gaussian condition in (16) can be weakened substantially.
Write \(D_{\rm tr}(\rho,\sigma)=\tfrac12\|\rho-\sigma\|_1\).
The following classical estimate makes the cost of an approximation
explicit, uniformly in block length. Common-mixture and inflated-energy
continuity methods have established antecedents \cite{winter,shirokov};
we give a direct proof with the particular channel cost needed here.

\begin{lemma}[Information continuity with an explicit cost]
\label{lem:tv}
Let \(W\) be a memoryless mixture of complex Gaussian channels with
\(|a|^2\leq K\) and component variance at least \(\ell>0\).
If laws \(Q,G\) on \(\C^m\) both have total second moment at most
\(mP\), and \(\operatorname{TV}(Q,G)\leq d\leq1\), then
\[
 |I_Q(W^{\otimes m})-I_G(W^{\otimes m})|
 \leq h_2(d)+md\log\left(1+\frac{KP}{\ell d}\right).       \tag{19}
\]
The product notation refers to independent branches of the same mixture.
The correction at \(d=0\) is defined as zero. An upper noise bound is
not required for this lemma. The estimate is uniform in the mixing law.
\end{lemma}
\begin{proof}
For any input with coordinate costs \(c_i\), reveal the independent
branch vector \(Z^m\) only to obtain an information upper bound.
Conditioned on a finite-variance branch vector, Gaussian maximum entropy,
subadditivity and Jensen give
\[
 0\leq I(X^m;Y^m)\leq I(X^m;Y^m,Z^m)
 =I(X^m;Y^m|Z^m)
 \leq\sum_i\log(1+Kc_i/\ell)
 \leq m\log(1+KP/\ell).
\]
These are relative-entropy chain rules. Conditional Gaussian smoothing
has bounded density and finite second moment, so the conditional
calculations are finite. After integrating the branch bound, mutual
information stays finite even when the mixed output entropy is infinite;
no subtraction of those mixed differential entropies is used.

For \(0<d<1\), a common positive submeasure of \(Q,G\) of mass \(1-d\)
exists: shrink their maximal common submeasure if their actual total
variation is less than \(d\). Normalize it to \(R\), and write
\(Q=(1-d)R+dS\), \(G=(1-d)R+dT\). Both exceptional laws have cost
at most \(mP/d\). A binary mixture label \(B\), followed by the channel,
gives exactly
\[
 I_Q=(1-d)I_R+dI_S+I_Q(B;Y^m),\qquad
 I_G=(1-d)I_R+dI_T+I_G(B;Y^m).
\]
The last terms lie in \([0,h_2(d)]\), and \(I_S,I_T\) lie in
\([0,m\log(1+KP/(\ell d))]\). Subtract their common parts to prove
(19), with one binary entropy. The endpoints follow from equality or
the preceding cost bound. This common-submeasure proof is essential:
the right side is not globally monotone in \(d\). A uniform bound on
differences also bounds differences of the infima over a channel family;
it does not interchange an input supremum and a channel infimum.
\end{proof}

For the matched-reset channel, let \(\bar\rho_i\) denote the average
input marginal of mode \(i\). Finite photon energy ensures finite first
moments and symmetrized quadrature covariance. The uncertainty inequality
\(V_i+i\Omega/2\geq0\) defines its unique moment-matched Gaussian
state \(\sigma_i\), with exactly the same energy
\(e_i=\Tr\bar\rho_i n\). This comparison allows squeezing.
Set
\[
 d_i=D_{\rm tr}(\bar\rho_i,\sigma_i),\quad
 d=m^{-1}\sum_i d_i,\quad P_i=(e_i+1)/E,\quad
 c_0=\lambda\log(1+k/r)+4k/r.
\]

\begin{theorem}[Robust marginal-Gaussian converse]\label{thm:robust}
For \(E\geq2k/r\), any \(m\)-mode ensemble of average total photon
cost at most \(mE\), with arbitrary conditional states and arbitrary
correlations in its average, satisfies
\[
 \frac{\chi(J:N_{E,\theta}^{\otimes m})}{m}
 \leq H_G+\frac{c_0}{E}+h_2(d)
       +d\log\left(1+\frac{2k(1+1/E)}{rd}\right).          \tag{20}
\]
The last two terms vanish at \(d=0\) and tend to zero as \(d\to0\),
uniformly in \(m\). No bound \(e_i\leq E\) is imposed on individual
modes. The same conclusion holds for continuous ensembles.
\end{theorem}
\begin{proof}
The simulation and continuous-label data processing used for (14) give
\[
 \chi\leq\sum_i I_{Q_{\bar\rho_i,E}}(W_{E,-,\theta}).
\]
Here heterodyne contraction bounds the total variation to
\(Q_{\sigma_i,E}\) by \(d_i\), and both laws have second moment \(P_i\).
Lemma~\ref{lem:tv}, with noise floor \(r/2\), therefore bounds each
replacement by
\[
 h_2(d_i)+d_i\log(1+2kP_i/(rd_i)).
\]
The nonmaximal common part in its proof justifies using \(d_i\),
without a false monotonicity assertion.

The comparison Husimi marginal is real planar Gaussian. Noise sensitivity
costs at most \(4k/(rE)\) per mode, uniformly in its cost. The Gaussian
hull combines all those costs before concavity is applied, giving
\(H_G(1+1/E)\leq H_G+\lambda\log(1+k/r)/E\).
Finally \((u,p)\mapsto u\log(1+(2k/r)p/u)\), extended by zero
when \(u=0\), is jointly concave: it is the perspective of
\(\log(1+(2k/r)p)\). Jensen for this function and for \(h_2\), and
\(m^{-1}\sum_iP_i\leq1+1/E\), prove (20). All mutual informations
are finite by the cost estimate in Lemma~\ref{lem:tv}; output quantum
energy is finite in the stated model. Thus no entropy ambiguity is
hidden in the chain rules.
\end{proof}

Define a distinct, larger functional \(R_{\rm marg}(E)\) by replacing
the global Gaussianity condition in (16) by Gaussianity of every
one-mode average marginal. The joint average need not be Gaussian.
Theorem~\ref{thm:robust} at \(d=0\) and the same product-Gaussian
lower construction used for (17) imply
\[
 H_G-\frac2{rE}\leq R_{\rm marg}(E)
 \leq H_G+\frac{c_0}{E}.                                  \tag{21}
\]
This is a Holevo functional, not an asserted operational capacity.
For approximately Gaussian marginals with \(d\leq\epsilon_m\to0\),
continuity at zero of the correction in (20) suffices for its uniform
vanishing; global monotonicity of that correction is unnecessary.

\subsection{A precise typical-codebook class}

A discrete empirical distribution of amplitudes has total variation one
from a continuous Gaussian. We instead control the empirical quantum
states and their moments separately.
Fix \(E>0\), a finite list of real planar Gaussian laws
\(G_l=N_2(u_l,V_l)\), \(V_l\geq0\), with
\(p_l=|u_l|^2+\Tr V_l<\infty\), and a deterministic known schedule
\(l_m(i)\), \(1\leq i\leq m\). This schedule carries no message;
the list and \(E\) are fixed before \(m\to\infty\).
Let \(M_m\geq e^{Rm}\), \(R>0\). Independently over messages and
coordinates, draw \(X_{ji}\sim G_{l_m(i)}\) and use equiprobable
product coherent words \(\bigotimes_i|\sqrt E X_{ji}\rangle\).
Independence between experiments at different lengths is not required.
Set \(P_m=m^{-1}\sum_i p_{l_m(i)}\).

\begin{proposition}[Typical marginal Gaussianity and physical cost]
\label{prop:typical}
For the preceding experiment, almost surely
\[
 \max_{i\leq m}D_{\rm tr}(\widehat\rho_i,
                \widehat\sigma_i)\longrightarrow0,
\]
where \(\widehat\rho_i\) is the empirical average marginal and
\(\widehat\sigma_i\) its own moment-matched Gaussian.
If \(P_m\leq1-\eta\) eventually for a fixed \(\eta>0\), the physical
average cost is at most \(mE\) almost surely eventually.
If only \(P_m\leq1\), the common correction below gives exact feasible
cost, preserves the displayed convergence, and changes the average error
of any given channel/decoder by \(o(1)\), uniformly in the channel.
No existence of a reliable decoder is asserted by this proposition.
\end{proposition}
\begin{proof}
Let \(\tau_l=\int|\sqrt E x\rangle\langle\sqrt E x|\,dG_l(x)\)
and \(N_*=E\max_l p_l\). Projection \(\Pi_L\) onto photon numbers
\(0,\ldots,L\) has rank \(L+1\) and
\(\Tr\tau_l(I-\Pi_L)\leq N_* /(L+1)\).
The nonnormalized truncation estimate
\(\|\rho-\Pi_L\rho\Pi_L\|_1\leq2\sqrt{\Tr\rho(I-\Pi_L)}\)
and Jensen control both empirical and design tails in expectation.
For projected independent rank-one samples, Hilbert--Schmidt variance
is at most \(1/M\), so their mean trace fluctuation is at most
\(\sqrt{(L+1)/M}\). Consequently, with \(M=M_m\),
\[
 \E\|\widehat\rho_i-\tau_{l_m(i)}\|_1
 \leq4\sqrt{\frac{N_*}{L+1}}+\sqrt{\frac{L+1}{M}}.
\]
Choose \(L+1\) of order \(\sqrt M\). Markov at threshold
\(M^{-1/8}\), then a union bound over \(m\) modes, gives failure
probability at most \(C_E mM_m^{-1/8}\), summable in \(m\).
Borel--Cantelli proves uniform convergence to the design states.

That alone does not control moment matching. Separately put
\(\widehat u_i=M^{-1}\sum_jX_{ji}\),
\(\widehat T_i=M^{-1}\sum_jX_{ji}X_{ji}^{\mathsf T}\), and
\(\widehat V_i=\widehat T_i-\widehat u_i\widehat u_i^{\mathsf T}\).
The finite Gaussian list has uniformly finite fourth moments, hence
\[
 \begin{aligned}
 \E|\widehat u_i-u_l|^2&\leq C/M,\\
 \E\|\widehat T_i-(V_l+u_lu_l^{\mathsf T})\|_F^2&\leq C/M.
 \end{aligned}
\]
Chebyshev and a union bound at threshold \(M_m^{-1/8}\) give failure
bounds \(CmM_m^{-3/4}\), summable in \(m\). Almost surely, means and raw
second moments converge uniformly with error \(O(M_m^{-1/8})\).
Thus covariances converge uniformly, even at a singular design covariance.

The moment-matched quantum Gaussian is the coherent mixture with
displacement law \(N_2(\widehat u_i,\widehat V_i)\).
For \([q,p]=i\), its mean is \(\sqrt{2E}\widehat u_i\) and covariance
is \(I/2+2E\widehat V_i\). Coupling two Gaussian laws by the same
standard real Gaussian vector and using the coherent overlap gives
\[
 D_{\rm tr}(\tau(u,V),\tau(u',V'))
 \leq\sqrt E\bigl(|u-u'|^2+
             \|\sqrt V-\sqrt{V'}\|_F^2\bigr)^{1/2}.
\]
The positive matrix square root is continuous on positive semidefinite
matrices, including the boundary: uniformly approximate the scalar square
root by polynomials on a compact spectral interval. The empirical
covariances stay eventually in such an interval. The last display and
the triangle inequality prove convergence to the empirical moment-matched
states, rather than merely to the design states.

The empirical cost is
\(\widehat P_m=(mM_m)^{-1}\sum_{ij}|X_{ji}|^2
=m^{-1}\sum_i\Tr\widehat T_i\).
Its difference from \(P_m\) is \(O(M_m^{-1/8})\) almost surely.
Fixed slack therefore proves eventual feasibility. For \(P_m\leq1\)
without slack put
\[
 t_m=\min(1,\widehat P_m^{-1/2}),\qquad X'_{ji}=t_mX_{ji},
\]
with \(t_m=1\) at zero cost. This is a common correction fixed with
the codebook, not a message-bearing operation. Its cost is exactly
feasible and \(1-t_m=O(M_m^{-1/8})\) almost surely.
The rescaled empirical means/covariances and their Gaussian states
converge to the same design values. Moreover coherent overlap and Jensen
bound the average whole-block trace perturbation by
\[
 M_m^{-1}\sum_jD_{\rm tr}(\rho'_j,\rho_j)
 \leq |1-t_m|\sqrt{mE\widehat P_m}
 =O(\sqrt{mE}M_m^{-1/8})=o(1).
\]
Trace-distance contraction controls the change in average decoding error
for any channel and any fixed decoder, uniformly over channels.
This transfers reliability if it was already proved; it does not create
it. The corrected random book is no longer iid, and is described
explicitly as the iid design followed by this common adjustment.
\end{proof}

\begin{corollary}[Operational consequence]\label{cor:operational}
At fixed \(E\geq2k/r\), any reliable equiprobable-message sequence
with cost \(mE\) and average marginal distances \(d_m\to0\) obeys
\[
 \limsup_{m\to\infty}\frac{\log M_m}{m}
 \leq H_G+\frac{c_0}{E}.                                  \tag{22}
\]
This applies almost surely to the feasible typical books in
Proposition~\ref{prop:typical}, if they have vanishing average error.
It includes collective decoders and the unknown-fixed-phase problem.
For sufficiently large \(E\), an unrestricted non-Gaussian common code
from (9) has a rate strictly above the right side of (22).
\end{corollary}
\begin{proof}
For error \(\epsilon_m\), Holevo and Fano give
\[
 (1-\epsilon_m)\log M_m/m-h_2(\epsilon_m)/m\leq\chi_m/m.
\]
Use (20) and then take the block limit at fixed \(E\).
For uniform error over phases, fix any phase for the converse.
The non-Gaussian strict-cost rate used in (15) is at least
\(H_G+\delta-2/(rE)\); it exceeds (22) when
\(E>(c_0+2/r)/\delta\), with the earlier admissibility conditions.
Its common decoder is provided by (9), not by the random-book theorem.
\end{proof}

\subsection{Two necessary cautions}

Trace closeness alone does not control the moment-matched Gaussian,
even with uniformly bounded energy. Indeed
\(\rho_L=(1-1/L)|0\rangle\langle0|+(1/L)|L\rangle\langle L|\)
has energy one and trace distance \(1/L\) to vacuum, while its Gaussian
with the same moments is always \(\gamma_1\).
Thus \(D_{\rm tr}(\rho_L,\gamma_1)\to1/2\), not zero.
This is why the empirical moments above were controlled separately.
Also, for a saturated circular design \(\CN(0,1)\), the actual cost is
the average of \(mM_m\) independent unit-mean exponentials, so
\(\Pr(\widehat P_m\leq1)\to1/2\) by the central limit theorem.
Design cost in expectation does not imply typical feasibility without
slack or correction.

No theorem here covers every deterministic realization named Gaussian,
arbitrary postselection, arbitrary message weights, message-bearing
covariance schedules, or design lists growing with block length.
No trace approximation of the whole average to a product Gaussian is
claimed or needed. Equation (21) remains a functional statement;
(22) is a separate converse for a defined operational code class,
not a proof of its exact capacity or random-code achievability.


\section{Prior work, limits and verification context}

Gaussian measure--prepare structure, the single Gaussian classical
limit, data processing, Fano, heat-flow score methods, and type coding
are prior tools \cite{holevoEB,holevoReview,types,compound}.
Recent general Gaussian-encoding optimality results for single-mode
Gaussian channels are given by Adesso, Lami and Giovannetti
\cite{adesso2026}. Their Gaussian-channel hypothesis excludes the
generic mixtures studied here; a mixture of Gaussian components is not
in general a Gaussian channel. Their fixed-error energy-rescaling
formula is not imported into our compound average-error problem.
Uniform quantization using compactness and Dini-type arguments is also
explicit in earlier compound-channel work \cite{mitran,schaefer}.
The Mitran--Devroye--Tarokh operational model assumes that the channel
parameter is known to the receiver; its theorem is therefore not a
substitute for our decoder without that side information.
Our use of its quantization mechanism is separately justified above.
Classical on/off fading already motivates non-Gaussian input
optimization \cite{tu}. Non-Gaussian advantages in entanglement-assisted
quantum fading are also prior research \cite{catalano}.
Near-capacity coherent modulation for non-Gaussian attenuator environments
and additive classical noise is also established \cite{huber}.
Those bounds depend on the environment and are independent of input
energy at fixed channel. Our joint scaling and unknown-law comparison
must be distinguished from that earlier setting; coherent near-optimality
itself is not new. Our example therefore does not discover the broad
non-Gaussian phenomenon.
Common-mixture continuity methods and exceptional-energy penalties have
antecedents in Winter \cite{winter} and Shirokov \cite{shirokov}.
Winter's energy-constrained bounds use Gibbs hypotheses; Shirokov's
ensemble distances do not simply become distances between average
marginals. Neither attribution replaces our direct proofs of (19)--(22).
The contribution here is the quantitative common simulation, Gaussian-hull
separation, and robust marginal converse for the stated joint scaling,
including the defined typical books and their physical cost correction.
These tools and that consequence are distinguished from a general
Gaussian-protocol impossibility theorem or an exact Gaussian-code capacity.

The results do not cover within-block changes of the unknown law,
channel memory, unbounded gains, a vanishing normalized noise floor,
arbitrary nonclassical reset states, or entanglement assistance.
Expected photon cost is not a deterministic photon-number bound.
A directed primary-source search is not an exhaustive literature audit.
The proof does not identify the general classical fading optimizer.

\paragraph{AI and responsibility disclosure.}
The named author directed substantial OpenAI Codex assistance in
derivation and writing. Technical audits requested GPT-6 Astra xhigh,
with shared filesystem and reused context declared.
Runtime identity, effort, isolation and statistical independence are
not independently attested. No independent human review, Lean
formalization, certified worldwide originality or predicted assessment
score is claimed. Assessment reports must be read with their exact
artifact hashes, declared participation and limitations. The unscored technical audit716 examined the new continuity, marginal
and typical-book arguments, including the two false auxiliary inferences
and the nonmaximal-common-part completion. It is authoring assistance,
not an independent vote. The published v1 and its assessment remain
unchanged. A full assessment of this exact revised PDF has not yet been
obtained; earlier manuscript scores do not apply to this revision.

\begin{thebibliography}{99}
\bibitem{holevoEB} A.~S.~Holevo, Entanglement-breaking channels in infinite
dimensions, \href{https://arxiv.org/abs/0802.0235v1}{arXiv:0802.0235v1} (2008).
\bibitem{holevoReview} A.~S.~Holevo and V.~Giovannetti, Quantum channels
and their entropic characteristics,
\href{https://arxiv.org/abs/1202.6480v1}{arXiv:1202.6480v1} (2012),
especially Sections 3.4--3.5 and 4.3.
\bibitem{types} I.~Csisz\'ar, The method of types, IEEE Transactions on
Information Theory 44, 2505--2523 (1998),
\href{https://doi.org/10.1109/18.720546}{doi:10.1109/18.720546}.
\bibitem{compound} I.~Bjelakovi\'c and H.~Boche, Classical capacities of
averaged and compound quantum channels,
\href{https://arxiv.org/abs/0710.3027v2}{arXiv:0710.3027v2} (2009),
Section II for the classical finite-alphabet formulation.
\bibitem{mitran} P.~Mitran, N.~Devroye and V.~Tarokh, On compound
channels with side-information at the transmitter,
\href{https://ece.uwaterloo.ca/~pmitran/compound3.pdf}{author draft,
20 December 2005}, Lemmas 3--5. Published in IEEE Transactions on
Information Theory 52, 1745--1755 (2006),
\href{https://doi.org/10.1109/TIT.2006.871044}{doi:10.1109/TIT.2006.871044}.
The cited draft is not identified as the final published byte version.
\bibitem{schaefer} R.~F.~Schaefer and S.~Loyka, The secrecy capacity
of compound Gaussian MIMO wiretap channels,
\href{https://arxiv.org/abs/1509.02182v1}{arXiv:1509.02182v1} (2015),
Section III for common quantization; its secrecy model and Gaussian
optimizer are not imported here.
\bibitem{tu} Tu, Li and Dhillon, An information-theoretic framework for
vision-aided communication over memoryless on-off fading channels,
npj Wireless Technology 2, 80 (2026),
\href{https://doi.org/10.1038/s44459-026-00074-7}
{doi:10.1038/s44459-026-00074-7}.
\bibitem{catalano} G.~Catalano, M.~Fanizza, F.~A.~Mele,
G.~De Palma and V.~Giovannetti, Convex combinations of bosonic
pure-loss channels,
\href{https://arxiv.org/abs/2604.26874v2}{arXiv:2604.26874v2} (2026),
selected discussion of entanglement-assisted fading and non-Gaussian
inputs. Version 2, revised 30 September 2026.
\bibitem{gaussianReview} C.~Weedbrook, S.~Pirandola,
R.~Garc\'ia-Patr\'on, N.~J.~Cerf, T.~C.~Ralph, J.~H.~Shapiro
and S.~Lloyd, Gaussian quantum information,
\href{https://arxiv.org/abs/1110.3234v1}{arXiv:1110.3234v1} (2011),
Sections II.E and II.E.2 in the PDF (II.5 and II.5.2 in numeric HTML)
for Gaussian measurement outcomes and heterodyne.
\bibitem{huber} S.~Huber and R.~K\"onig, Coherent state coding approaches
the capacity of non-Gaussian bosonic channels,
\href{https://arxiv.org/abs/1710.06734v2}{arXiv:1710.06734v2} (2018),
Theorem 1. Journal of Physics A 51, 184001,
\href{https://doi.org/10.1088/1751-8121/aab7ff}{doi:10.1088/1751-8121/aab7ff}.
\bibitem{adesso2026} G.~Adesso, L.~Lami and V.~Giovannetti,
Classical capacity of single-mode bosonic Gaussian channels,
\href{https://arxiv.org/abs/2610.22308v1}{arXiv:2610.22308v1}
(15 October 2026), Section IV and Theorem 2 for the Gaussian-channel scope.
\bibitem{holevoMeasurement} A.~Holevo, On the classical capacity of
general quantum Gaussian measurement, Entropy 23, 377 (2021),
\href{https://doi.org/10.3390/e23030377}{doi:10.3390/e23030377},
Theorem 1 and Remark 2 for the distinction between Gaussian average
and Gaussian ensemble; its Gaussian-measurement optimization is not
imported here.
\bibitem{winter} A.~Winter, Tight uniform continuity bounds for quantum
entropies: conditional entropy, relative entropy distance and energy
constraints, \href{https://arxiv.org/abs/1507.07775v6}{arXiv:1507.07775v6}
(12 January 2016), Lemma 2 and the energy-constrained discussion,
Proposition 14 and Lemma 15.
\bibitem{shirokov} M.~E.~Shirokov, Uniform continuity bounds for
information characteristics of quantum channels depending on input
dimension and on input energy,
\href{https://arxiv.org/abs/1610.08870v5}{arXiv:1610.08870v5}
(3 November 2018), Section 6.2, Propositions 7--8 for ensemble-distance
continuity; those hypotheses are not identified with marginal closeness.
\end{thebibliography}
\end{document}
